\section{Introduction}

\begin{figure}[t]
\centering
\includegraphics[width=0.72\linewidth]{fig_teaser.png}
\caption{From a handful of RGB views to an open-vocabulary 3D scene graph. \textbf{Left:} three input
frames (TUM \texttt{freiburg3\_long\_office\_household}). \textbf{Right:} frozen-model detections lifted
into VGGT world-point geometry and fused across views into labeled 3D object nodes (colored markers),
with all heavy models frozen and no per-scene optimization. Best viewed in color.}
\label{fig:teaser}
\end{figure}


Reconstructing a scene and explaining a scene differ. Feed-forward geometry models
recover camera motion, depth, and point maps from a few
images~\cite{wang2024dust3r,leroy2024mast3r,wang2025vggt}, and 2D foundation models give
open-vocabulary detections, masks, and
descriptors~\cite{minderer2023owlv2,radford2021clip,kirillov2023segment,oquab2023dinov2}. Many
applications instead need object-level structure: which pixels across views belong to the same
object, where it sits in 3D, and how it relates to others, so a robot can act on
``the monitor is on the desk'' rather than a raw point cloud. These components do not provide
this alone, and the gap is widest for a handful of casual RGB views rather than a
calibrated RGB-D stream or prebuilt map.

We build the scene graph by fusing lifted 2D detections, keeping every heavy model frozen with no
per-scene optimization. VGGT world-point maps act as a 2D-to-3D lookup, an open-vocabulary detector
proposes labeled object boxes, and we lift and fuse these across views into an object--relation graph
(Figure~\ref{fig:teaser}). Fixed perception isolates the graph-building stage for ablation and
places us in a lower-input regime than open-vocabulary 3D systems assuming posed RGB-D video
or a dense map~\cite{gu2024conceptgraphs,jatavallabhula2023conceptfusion}.

A central lesson of this work concerns evaluation. An earlier version graded
predicted graphs against a reference seeded from the model's own dense-view output, which is circular: a
method reproducing its ten-view graph scores $1.0$ at the reference view by construction, rewarding
agreement with the model rather than the scene. Rebuilding the
reference from the frames themselves changes two conclusions. First, a mask-plus-CLIP front-end recovers
almost no real objects; on a desk scene it returns ``floor'', ``curtain'', and ``bed''
instead of the present monitors and keyboard, whereas a genuine open-vocabulary detector
(OWLv2) reads the actual objects. Second, an uncertainty-aware fusion gate that looked promising
under the old protocol gives no improvement once labels are real, across a full weight sweep.
With honest labels, the dominant error is over-counting: one physical object survives
fusion as several same-label nodes (Figure~\ref{fig:dedup3d}). A one-line merge of same-label nodes
with overlapping 3D boxes removes most of it, the only variant whose accuracy improves with views.

We test these on 28 TUM RGB-D sequences~\cite{sturm2012tumrgbd} at $3/5/8/10$ views, against a
reference authored from the frames (5 scenes human-verified, 23 drafted in two passes). Duplicate-aware
fusion improves object-label F1 over the graph-fusion baseline by $0.06$ to $0.12$ and wins on $22$ to
$27$ of the $28$ scenes ($p<0.001$), while the uncertainty gate is consistently worse. We make four
contributions. \textbf{(1)}~A detector-grounded open-vocabulary scene-graph system for sparse views that
keeps OWLv2 and VGGT frozen with no per-scene optimization, recovering real
objects a mask-plus-CLIP front-end misses. \textbf{(2)}~Duplicate-aware cross-view fusion, a single
post-fusion merge of same-label nodes with overlapping 3D boxes that targets over-counting and
is the only variant whose accuracy improves with more views. \textbf{(3)}~An independent,
de-circularized benchmark on TUM RGB-D, showing how a prediction-seeded reference inflates scores
while an independently authored one removes the effect. \textbf{(4)}~A negative result on
uncertainty-aware fusion, with analysis: the gate raises recall but
over-splits, costing more precision than it gains, so merging helps.